\noindent\makebox[0pt][l]{\fontsize{13}{16}\selectfont 418}\hfill{\small\textbf{POSTHUMOUS PAPERS.}}\hfill\null\par\medskip

\begin{center}
A historical report\\
on the surveys partly carried out and partly directed by Court Councillor \textsc{Gauss}\\
in the Kingdom of Hanover.
\end{center}

\par\noindent\hspace*{\parindent}The various surveys of which I am to give a brief historical account here are closely connected, though their purposes and character differ. The clearest way to present them will be to arrange them in the chronological order of the commissions under which I undertook them.

\par\medskip
\begin{center}
I.\quad The Hanover arc measurement.
\end{center}

\par\noindent\hspace*{\parindent}By a rescript of 30 June 1820 I was commissioned to carry out an arc measurement through the Kingdom of Hanover, as an extension or continuation of a similar undertaking begun shortly before in the Danish territories. By that term one means those operations, partly astronomical and partly trigonometrical, by which the magnitude of a degree of the meridian is determined in a known unit of length (foot, toise, etc.). For this purpose one chooses two points on the same meridian that are sufficiently far apart, determines the length of the arc between them in feet, etc., by a network of triangles carried between them, and finds the number of degrees, minutes and seconds corresponding to the same arc by astronomical observations at the endpoints. From this one then infers the length of one degree.

\par\noindent\hspace*{\parindent}It follows that an arc measurement, \textls[120]{as such}, has a purely scientific aim, to contribute to mathematical knowledge of the proportions of the terrestrial spheroid. To meet the demands of the century, such knowledge requires much greater precision and many more arc measurements than formerly had to suffice. It also follows that an arc measurement consists of two very different parts, one trigonometrical and the other astronomical. Without going into further detail here, I will only remark that I brought the trigonometrical part (comprising 25 principal triangles from Göttingen to Hamburg) in 1821, 1822, 1823, and the astronomical part in 1827,

\par\clearpage
\noindent\null\hfill{\small\textbf{ON THE HANOVER TRIANGULATION.}}\hfill\makebox[0pt][r]{\fontsize{13}{16}\selectfont 419}\par\medskip

\noindent to full completion, and that the latter and the final results were made known in a work I published in 1828, »Bestimmung des Breitenunterschiedes zwischen den Sternwarten von Göttingen und Altona durch Beobachtungen am \textsc{Ramsden}schen Zenithsector«. The Göttingen and Altona observatories, which by a singular play of chance lie on precisely the same meridian, are the endpoints of the measured meridian arc, whose total curvature is somewhat more than two degrees.

\begin{center}
II.\quad Trigonometrical connection of the arc-measurement triangles\\
with the triangles of the royal Dutch surveys.
\end{center}

\par\noindent\hspace*{\parindent}The immediate scientific purpose of this extension of the trigonometrical work, assigned to me by a rescript of 8 March 1824, was to connect the Hanover arc measurement with the French and English ones. These latter, as is well known, are connected to one another, and the carefully measured triangles in the Netherlands join the French triangles. During the period in which East Frisia was united with Holland, these triangles had been carried as far as the eastern boundary of that principality, and farther south as far as Bentheim. I thus had the choice of two routes by which to achieve the intended purpose, either from the northernmost triangles of the arc measurement through Bremen to Jever, or from the middle triangles through Westphalia to Bentheim. The former route was preferred, among other reasons, because in that way a connection with the North Sea could also be obtained, and with it the determination of the absolute elevations of all the triangulation stations above sea level. I carried out this triangulation in 1824 and 1825, bringing the total number of principal triangles to 38[\textsuperscript{*})].

\par\medskip
\par\noindent\hspace*{\parindent}I have a few remarks to add about these trigonometrical measurements of 1821---1825, as they differ in several respects from similar work previously carried out in other countries.

\par\smallskip\noindent\rule{.24\linewidth}{.4pt}\par\smallskip
{\footnotesize\noindent\hspace*{\parindent}[\textsuperscript{*}) The number of independent triangles of the arc measurement between Göttingen and Hamburg is 21, and the total number of independent triangles of the arc measurement and its continuation to Jever (excluding the triangle Hohehagen-Brocken-Inselsberg) is 42. Cf. p. 297.]\par}

\par\smallskip\noindent\hfill{\footnotesize 53\textsuperscript{*}}\hspace*{.12\linewidth}\par

\par\clearpage
\noindent\makebox[0pt][l]{\fontsize{13}{16}\selectfont 420}\hfill{\small\textbf{POSTHUMOUS PAPERS.}}\hfill\null\par\medskip

\par\noindent\hspace*{\parindent}1) The use of the \textls[120]{heliotropes} I first introduced made special signal towers unnecessary, made it possible to use triangles of a size previously impracticable, and allowed a precision of measurement unattainable without that aid.

\par\noindent\hspace*{\parindent}2) In carrying out the measurements I did not restrict myself to what was necessary for the immediate purpose (which, as already noted, was primarily scientific). I also sought to make the measurements as productive for the geography of the country as possible without impairing that immediate purpose. Most of the towers in the area covered by the triangulation stations have therefore been located with great precision. These include the towers of the towns lying along the chain of triangles, such as Göttingen, Hildesheim, Wolfenbüttel, Braunschweig, Hannover, Celle, Lüneburg, Harburg, Hamburg, Buxtehude, Stade, Verden, Bremen, Oldenburg, Varel, Jever, and those of many hundreds of smaller localities, as the overview maps accompanying my reports on each year's work show. Fixed control points and a basis have thus been obtained for all future detailed surveys in the districts concerned, giving these determinations lasting value.

\par\noindent\hspace*{\parindent}3) Finally, by a method of my own I have put these results into a form that offers substantial advantages for the further use just described.

\begin{center}
III.\quad Trigonometrical survey of the districts\\
not covered by the measurements of 1821---1825.
\end{center}

\par\noindent\hspace*{\parindent}The work carried out in 1821---1825, though not primarily intended for this purpose, included an actual trigonometrical survey of a very considerable part of the Kingdom of Hanover. The generally acknowledged advantages of an accurate trigonometrical survey of the country made it desirable to go farther. By a rescript of 28 April 1828[\textsuperscript{*})] I was commissioned to direct the further extension of the trigonometrical survey to all districts not yet covered by the earlier work.

\par\smallskip\noindent\rule{.24\linewidth}{.4pt}\par\smallskip
{\footnotesize\noindent\hspace*{\parindent}[\textsuperscript{*}) The rescript in the Gauss archive is dated 14 April.]\par}

\par\clearpage
\noindent\null\hfill{\small\textbf{ON THE HANOVER TRIANGULATION.}}\hfill\makebox[0pt][r]{\fontsize{13}{16}\selectfont 421}\par\medskip

\par\noindent\hspace*{\parindent}At the same time a detailed survey was ordered, to cover all the districts that had not yet been included in the plane-table survey of the former Electorate of Hanover conducted in the eighties of the previous century. Colonel \textsc{Prott} was charged with direct supervision of this detailed survey. The work is directly connected with the trigonometrical survey in that the results of the latter provide the basis and fixed control points for the plane-table work. The ease and reliability that the plane-table work gains in this way had already been demonstrated a year earlier in some preliminary detailed surveys in the Hildesheim region, even though the work then was carried out by officers who had previously been unfamiliar with the method.

\par\noindent\hspace*{\parindent}What I have to say here about the trigonometrical survey placed under my direction concerns the personnel engaged in it, the extent of the measurements completed so far, and finally the measurements still outstanding.

\begin{center}
A.\quad The personnel.
\end{center}

\par\noindent\hspace*{\parindent}A trigonometrical survey requires two entirely different kinds of work, taking the measurements at the relevant sites in the field, and processing them into results by combination and calculation at the desk. I have hitherto undertaken the second part of the work entirely myself and entrusted the first to the artillery officers who served as my assistants in 1821---1825. They had thus had the opportunity to become familiar both with the handling of the instruments and with the spirit and distinctive features of my procedure. These officers were Captain \textsc{Müller}, then First Lieutenant \textsc{Hartmann}, and my eldest son, now First Lieutenant in the artillery regiment. Lieutenant \textsc{Hartmann}, who in 1831 left military service for a post at the higher trade school in Hannover with the titular rank of captain, died in 1834.

\par\clearpage
\noindent\makebox[0pt][l]{\fontsize{13}{16}\selectfont 422}\hfill{\small\textbf{POSTHUMOUS PAPERS.}}\hfill\null\par\medskip

\begin{center}
B.\quad Measurements of the principal triangles completed so far.
\end{center}

\par\noindent\hspace*{\parindent}Had the three officers named been able to devote themselves exclusively and without interruption to this work throughout the suitable season each year, it would long since have been completed. But various impediments, due partly to the officers' duties and partly to other external circumstances, meant that in some years they could work on it only for a shorter period and in others could do nothing at all. In addition to the already mentioned death of Captain \textsc{Hartmann}, the ongoing detailed survey also continued to require special preparatory measurements, especially in districts where the church towers are more widely scattered. For these tasks only one or another of those officers could be used, and his time was then taken away from the main undertaking. The following is a summary of what has so far been accomplished in each year.

{\hangindent=.075\linewidth\hangafter=1\noindent\makebox[.075\linewidth][l]{1828}The work began in late summer. Captain \textsc{Müller} and Lieutenant \textsc{Gauss} completed most of the trigonometrical survey of Eichsfeld, while Lieutenant \textsc{Hartmann} took preparatory measurements for the detailed survey in the district of Hunnesrück and part of the Hildesheim region.\par}

{\hangindent=.075\linewidth\hangafter=1\noindent\makebox[.075\linewidth][l]{1829.}\enspace Lieutenant \textsc{Gauss} first made a few supplementary measurements in Eichsfeld. He devoted most of the year, however, to triangulation in Westphalia, where he carried out a system of triangles from the Weser to Bentheim. Captain \textsc{Müller} had been able to participate only briefly in the same undertaking, as he was called away from the survey by other assignments shortly after it began. Lieutenant \textsc{Hartmann}'s work in that year consisted of preparatory measurements for the detailed survey, first in the Hildesheim region, then in the districts of Uchte, Freudenberg and Auburg.\par}

{\hangindent=.075\linewidth\hangafter=1\noindent\makebox[.075\linewidth][l]{1830.}\enspace Trigonometrical survey of the eastern part of the Lüneburg region by Captain \textsc{Müller} and Lieutenant \textsc{Gauss}. Extension to East Frisia of the chain of triangles measured through Westphalia the previous year, by Lieutenant \textsc{Hartmann}.\par}

{\hangindent=.075\linewidth\hangafter=1\noindent\makebox[.075\linewidth][l]{1831}Only a short time could be devoted to the work. Lieutenant \textsc{Gauss}\par}

\par\clearpage
\noindent\null\hfill{\small\textbf{ON THE HANOVER TRIANGULATION.}}\hfill\makebox[0pt][r]{\fontsize{13}{16}\selectfont 423}\par\medskip

{\leftskip=.075\linewidth\noindent completed the measurements in the Lüneburg region, Lieutenant \textsc{Hartmann} those in East Frisia. Captain \textsc{Müller} took no part at all.\par}

{\hangindent=.075\linewidth\hangafter=1\noindent\makebox[.075\linewidth][l]{1832}The measurements were suspended altogether.\par}

{\hangindent=.075\linewidth\hangafter=1\noindent\makebox[.075\linewidth][l]{1833.}\enspace Triangulation of the district along the Weser from the vicinity of Nienburg to Holzminden by Captain \textsc{Müller} and Lieutenant \textsc{Gauss}. Triangulation of the Harz by Captain \textsc{Hartmann}.\par}

{\hangindent=.075\linewidth\hangafter=1\noindent\makebox[.075\linewidth][l]{1834}Only Lieutenant \textsc{Gauss} could devote even a very short time to the work, partly to make preparatory measurements for the detailed survey in the Osnabrück region and partly to supplement the previous year's measurements on the Weser.\par}

{\hangindent=.075\linewidth\hangafter=1\noindent\makebox[.075\linewidth][l]{1835.}\enspace This year again only Lieutenant \textsc{Gauss} could be spared for a short time to take further preparatory measurements in the Osnabrück region, this time with the assistance of Dr. \textsc{Goldschmidt}, observer at the Göttingen observatory.\par}

{\leftskip=.075\linewidth\noindent In the two following years only Captain \textsc{Müller} could devote himself to the measurements, namely\par}

{\hangindent=.075\linewidth\hangafter=1\noindent\makebox[.075\linewidth][l]{1836}to the triangulation of the district on the upper Weser between Uslar, Göttingen and Münden,\par}

{\hangindent=.075\linewidth\hangafter=1\noindent\makebox[.075\linewidth][l]{1837}to further preparations for the detailed survey in the Osnabrück region.\par}

{\leftskip=.075\linewidth\noindent In addition, Captain \textsc{Müller} carried out a preliminary reconnaissance of the Aller district that year for a future trigonometrical survey of it.\par}

\begin{center}
C.\quad Measurements of principal triangles still outstanding.
\end{center}

\par\noindent\hspace*{\parindent}The trigonometrical surveys enumerated so far are all connected with one another. To cover the \textls[120]{whole} kingdom and form a whole complete in its essentials, all that is still needed is the trigonometrical survey of two larger districts, namely, first, the area to the right and left of the Aller, and second, the northern part of the Bremen region. As already mentioned, a preliminary reconnaissance has been made in the first of these districts.

\par\noindent\hspace*{\parindent}In work of this kind it is impossible to determine in advance with any precision the time and expense required, since the rate of progress depends on many circumstances, some

\par\clearpage
\noindent\makebox[0pt][l]{\fontsize{13}{16}\selectfont 424}\hfill{\small\textbf{POSTHUMOUS PAPERS.}}\hfill\null\par\medskip

\noindent not known exactly in advance and some accidental (for example particular difficulties of the terrain or the state of the weather). It may nevertheless be assumed with some probability that if Captain \textsc{Müller} is to devote himself to these measurements alone, he could complete them entirely or for the most part in two summers.

\par\medskip
\par\noindent\hspace*{\parindent}The final results proper of the trigonometrical measurements themselves are the \textls[120]{numbers} that fix with the greatest precision the positions of the objects determined (mostly church towers) and therefore retain lasting value. Once the work is fully completed, they can be arranged in order and deposited in several copies for greater security, one of them perhaps at the observatory. Whether they should also be published in print will then depend on a decision by higher authority. I cannot leave unmentioned, however, that \textsc{Papen}'s map of the kingdom, which equals the most accomplished works of its kind in every respect and of which 22 sheets (thus one third of the whole) have already appeared, owes its accuracy to those results. All the results obtained so far are available for this purpose to Lieutenant \textsc{Gauss}, who has undertaken to graduate that map and plot all the fundamental positions on it precisely.

\par\medskip
\par\noindent\hspace*{\parindent}In conclusion I wish to add something about the preparatory measurements for the detailed survey. Their date of completion cannot be determined in advance at all, because by the nature of the work they necessarily have to continue almost as long as the detailed survey itself. For in districts where the church towers determined by the principal triangulation are too far apart and too few in number to provide enough fixed control points for each plane-table sheet, special signal stakes must be set up and their positions determined by smaller triangles. But there is no way to protect such signal stakes so securely that one can count with certainty on their remaining in place for many years. Although the loss of a single signal stake here or there might not render the work useless,

\par\clearpage
\noindent\null\hfill{\small\textbf{ON THE HANOVER TRIANGULATION.}}\hfill\makebox[0pt][r]{\fontsize{13}{16}\selectfont 425}\par\medskip

\noindent the disappearance of several before their use in the plane-table survey would mean the loss of all the work spent in determining their positions. Such preparatory operations must therefore be undertaken at most one or a few years before the district in question comes up for detailed survey.

\par\noindent\hspace*{\parindent}The detailed survey itself, incidentally, is complete in the Hildesheim region, Eichsfeld, the district of Hunnesrück and the districts of Freudenberg, Uchte and Auburg, and is already fairly far advanced in the Principality of Osnabrück. Only Colonel \textsc{Prott}, however, can give a more exact and complete account of it, as the work is within his remit.

\par\noindent\hspace*{\parindent}Göttingen, 8 February 1838.

\par\medskip
\begin{center}
[From the report of 5 July 1840 ȟber die trigonometrischen\\
Vermessungen« in 1839.]
\end{center}

\par\noindent\hspace*{\parindent}\mbox{.\ .\ .\ .\ .\ .}\quad In assessing the systems of triangles one must bear in mind that the original intention was not a general survey of the country, but at first only an arc measurement from Göttingen to Holstein, and then, initially, an extension of the system of triangles to join the \textsc{Krayenhoff} triangles as far as East Frisia. In keeping with \textls[120]{these} purposes, the triangles from Inselsberg to Jever that I myself measured in 1821---1825 were \mbox{chosen.\enspace\mbox{.\ .\ .\ .}}\quad The remaining triangles, added later, appear as branches of those principal triangles. As a result of this manner of formation, the system as a whole does not everywhere appear as a rounded whole cast in one piece, to the extent that it would have done if such a consideration had had to be taken into account from the outset. But the actual purpose, namely the precise fixing of the principal points in the country as a whole that are suited to this use, has been achieved just as \mbox{well.\enspace\mbox{.\ .\ .\ .\ .\ .}}

\par\medskip
\begin{center}
[From a report, December 1844, »über die im Jahre 1844 ausgeführten\\
trigonometrischen Messungen.«]
\end{center}

\par\noindent\hspace*{\parindent}\mbox{.\ .\ .\ .\ .\ .}\quad I take the liberty of adding a few words about all the measurements of the past years.

\par\smallskip\noindent\hspace*{.03\linewidth}{\footnotesize IX.}\hfill{\footnotesize 54}\hspace*{.12\linewidth}\par

\par\clearpage
\noindent\makebox[0pt][l]{\fontsize{13}{16}\selectfont 426}\hfill{\small\textbf{POSTHUMOUS PAPERS.}}\hfill\null\par\medskip

\par\noindent\hspace*{\parindent}\mbox{.\ .\ .\ .\ .\ .}\quad Each year after the measurements have been processed, the results [i.e. the coordinates] have been entered in lists. There are sixteen such partial lists, containing altogether somewhat more than 3000 determinations. The number of points itself may, however, be smaller by about one seventh, since many points that could be determined more precisely or reliably in a later year when new data became available appear in more than one list. There can be few church towers in the whole kingdom left undetermined.

\par\noindent\hspace*{\parindent}There is no need to explain further here that these lists have formed the basis of all detailed surveys made over the past 16 or 17 years and of \textsc{Papen}'s maps. But it is of great importance that these numbers, which retain their value as long as the objects exist, cannot be lost. The lists mentioned are kept at the observatory. Lieutenant \textsc{Gauss}, who entered all the fundamental points on \textsc{Papen}'s maps, has copies of them as well. For greater security and ease of use I have now begun to combine the partial lists into one, which will accordingly contain about 2600 points.

\par\noindent\hspace*{\parindent}It might be considered advisable later to publish this list or an extract from it in print. For the present, however, I would still consider this premature for several reasons.

\par\noindent\hspace*{\parindent}First, because a \textls[120]{scientifically} adequate exposition of the significance and full range of uses of these numbers can be given only gradually, in connection with the development of my own mathematical theories, which I intend to present in a series of separate papers (perhaps three or four). The first of these has already been printed[\textsuperscript{*})] as part of the forthcoming volume of the Denkschriften of the local Society of Sciences and has also been put on sale separately. I shall follow it with the others in turn as soon as possible.

\par\noindent\hspace*{\parindent}Second, because the numbers in the list, though sufficiently, indeed more than sufficiently precise for any practical use, still allow some refinement and correction when the various measurements are merged into one system that meets the strictest

\par\smallskip\noindent\rule{.24\linewidth}{.4pt}\par\smallskip
{\footnotesize\noindent\hspace*{\parindent}[\textsuperscript{*}) Untersuchungen über Gegenstände der höhern Geodäsie. Erste Abhandlung. Presented to the Royal Society on 23 October 1843. Volume IV, pp. 259/300.]\par}

\par\clearpage
\noindent\null\hfill{\small\textbf{ON THE HANOVER TRIANGULATION.}}\hfill\makebox[0pt][r]{\fontsize{13}{16}\selectfont 427}\par\medskip

\noindent theoretical requirements, even if the changes are in themselves very slight.

\par\noindent\hspace*{\parindent}Third, because I would like to give the position of the most important points, particularly the church towers in towns, in another form besides their coordinates, namely by geographic latitude and longitude. This transformation always requires considerable time and so can be carried out only \mbox{gradually.\quad\mbox{.\ .\ .\ .\ .\ .}}

\par\medskip
\begin{center}
[3.]\\
{\large Principal triangulation stations of the Hanover surveys.}
\end{center}

\par\smallskip
{\centering
\setlength{\tabcolsep}{0pt}
\newcommand{\GaussIXCoordinate}[7]{%
  \makebox[\linewidth][c]{%
    \makebox[1.8em][r]{\(#1\)\makebox[0pt][l]{\(\vphantom{0}#5\)}}%
    \hspace{.2em}\makebox[1.35em][r]{\(#2\)\makebox[0pt][l]{\(\vphantom{0}#6\)}}%
    \hspace{.2em}\makebox[1.55em][r]{\(#3\)\makebox[0pt][l]{\(\vphantom{0}#7\)}}%
    \( {,}#4 \)%
  }%
}
\begin{tabular}{@{\vrule width .8pt}p{.34\linewidth}|p{.23\linewidth}|p{.27\linewidth}@{\vrule width .8pt}}
\noalign{\hrule height .8pt}
\rule[-1.0\baselineskip]{0pt}{2.6\baselineskip}\makebox[\linewidth][c]{Triangulation station} &
\makebox[\linewidth][c]{Latitude} &
\makebox[\linewidth][c]{Long.\ from Gött.\ merid.} \\
\hline
\multicolumn{1}{@{\vrule width .8pt}c|}{\rule{0pt}{.45\baselineskip}} &
& \\
\hspace*{.6em}Inselsberg, Hess. tri. pt. & \GaussIXCoordinate{50}{51}{8}{619}{^{0}}{^{\prime}}{^{\prime\prime}} & \GaussIXCoordinate{-0}{31}{24}{311}{^{0}}{^{\prime}}{^{\prime\prime}} \\
\hspace*{.6em}Hohehagen & \GaussIXCoordinate{51}{28}{31}{234}{}{}{} & \GaussIXCoordinate{+0}{10}{45}{143}{}{}{} \\
\hspace*{.6em}Göttingen Observatory & \GaussIXCoordinate{51}{31}{48}{028}{}{}{} & \makebox[\linewidth][c]{\(0\)} \\
\hspace*{.6em}Meridian mark & \GaussIXCoordinate{51}{34}{30}{309}{}{}{} & \GaussIXCoordinate{-0}{0}{0}{007}{}{}{} \\
\hspace*{.6em}Brocken & \GaussIXCoordinate{51}{48}{1}{849}{}{}{} & \GaussIXCoordinate{-0}{40}{23}{071}{}{}{} \\
\hspace*{.6em}Hils & \GaussIXCoordinate{51}{53}{53}{000}{}{}{} & \GaussIXCoordinate{+0}{6}{41}{152}{}{}{} \\
\hspace*{.6em}Lichtenberg & \GaussIXCoordinate{52}{7}{22}{020}{}{}{} & \GaussIXCoordinate{-0}{20}{33}{348}{}{}{} \\
\hspace*{.6em}Deister & \GaussIXCoordinate{52}{14}{5}{778}{}{}{} & \GaussIXCoordinate{+0}{20}{35}{702}{}{}{} \\
\hspace*{.6em}Garssen & \GaussIXCoordinate{52}{39}{40}{397}{}{}{} & \GaussIXCoordinate{-0}{12}{19}{154}{}{}{} \\
\hspace*{.6em}Scharnhorst & \GaussIXCoordinate{52}{43}{43}{000}{}{}{} & \GaussIXCoordinate{-0}{19}{1}{620}{}{}{} \\
\hspace*{.6em}Breithorn & \GaussIXCoordinate{52}{49}{45}{985}{}{}{} & \GaussIXCoordinate{-0}{19}{11}{230}{}{}{} \\
\hspace*{.6em}Falkenberg & \GaussIXCoordinate{52}{50}{52}{481}{}{}{} & \GaussIXCoordinate{+0}{4}{34}{858}{}{}{} \\
\hspace*{.6em}Hauselberg & \GaussIXCoordinate{52}{51}{36}{513}{}{}{} & \GaussIXCoordinate{-0}{14}{28}{226}{}{}{} \\
\hspace*{.6em}Elmhorst & \GaussIXCoordinate{52}{59}{42}{886}{}{}{} & \GaussIXCoordinate{+0}{18}{24}{543}{}{}{} \\
\hspace*{.6em}Steinberg & \GaussIXCoordinate{52}{59}{54}{987}{}{}{} & \GaussIXCoordinate{+0}{40}{7}{347}{}{}{} \\
\hspace*{.6em}Bottel & \GaussIXCoordinate{53}{2}{22}{919}{}{}{} & \GaussIXCoordinate{+0}{39}{22}{916}{}{}{} \\
\hspace*{.6em}Wulfsode & \GaussIXCoordinate{53}{4}{15}{113}{}{}{} & \GaussIXCoordinate{-0}{17}{48}{837}{}{}{} \\
\hspace*{.6em}Bremen, Ansgar. Tower & \GaussIXCoordinate{53}{4}{48}{867}{}{}{} & \GaussIXCoordinate{-1}{8}{22}{781}{}{}{} \\
\multicolumn{1}{@{\vrule width .8pt}c|}{\rule{0pt}{.45\baselineskip}} &
& \\
\end{tabular}\par
}

\par\smallskip\noindent\hfill{\footnotesize 54\textsuperscript{*}}\hspace*{.12\linewidth}\par
